\newcommand{\ee}{\mathrm{e}}
\newcommand{\ii}{\mathrm{i}}
\heading{数学物理方法（上）-第10次作业}

\begin{question}
判断下列函数在 \(\infty\) 点处的性质，其中 \(n\in\mathbb{Z}\)：

\begin{enumerate}
\item \(z^n\);
\item \(\sin z^n,\; n\in\mathbb{Z}\);
\item \(\ee^{z^n}\);
\item \(\cos\sqrt{z^n}\);
\item \(\sin\sqrt{z^n}\);
\item \(\displaystyle\sqrt{\frac{z-1}{z-2}}\);
\item \(\sqrt{z-1}\);
\item \(\sqrt{(z-1)(z-2)}\).
\end{enumerate}

\begin{solution}
1. 做代换 \(t=1/z\)，有 \(z^n=t^{-n}\)，显然函数在 \(\infty\) 处为 \(n\) 阶极点。

2. 当 \(z\to\infty\) 时，\(z^n\to\infty\)，而函数 \(\sin z^n\) 极限不存在，故函数在 \(\infty\) 处为本性奇点。

3. 当 \(z\to\infty\) 时，\(z^n\to\infty\)，而函数 \(\exp(z^n)\) 极限不存在，故函数在 \(\infty\) 处为本性奇点。

4. 当 \(z\to\infty\) 时，\(\sqrt{z^n}\to\infty\)，而函数 \(\cos(\sqrt{z^n})\) 极限不存在，故函数在 \(\infty\) 处为本性奇点。

5. 当 \(z\to\infty\) 时，\(\sqrt{z^n}\to\infty\)，而函数 \(\sin(\sqrt{z^n})\) 极限不存在，故函数在 \(\infty\) 处为本性奇点。

6. 当 \(z\to\infty\) 时，\(\displaystyle\frac{z-1}{z-2}\to 1\)，因此有
   \[
   \lim_{z\to\infty} \sqrt{\frac{z-1}{z-2}} = \pm 1
   \]
   （取决于分支），故函数在 \(\infty\) 处为可去奇点。

7. 当 \(z\to\infty\) 时，\(\sqrt{z-1}\to\infty\)，故函数在 \(\infty\) 处为极点。做代换 \(t=1/z\)，
   \[
   \sqrt{\frac{1}{t}-1} = t^{-1/2}\sqrt{1-t},
   \]
   因此有 \(t=0\) 为 \(1/2\) 阶奇点，即 \(\infty\) 为 \(1/2\) 阶奇点。

8. 做代换 \(t=1/z\)，有
   \[
   \sqrt{(z-1)(z-2)} = \frac{1}{t}\sqrt{(1-t)(1-2t)},
   \]
   故函数在 \(t=0\) 处为 \(1\) 阶奇点，即 \(\infty\) 为 \(1\) 阶奇点。
\end{solution}
\end{question}

\begin{question}
已知：
\[
f(z)=\sum_{n=0}^\infty z^{n!}=z+z+z^2+z^6+z^{24}+\cdots,\qquad |z|<1.
\]

1. 证明：\(z=1\) 是 \(f(z)\) 的奇点。
2. 证明：\(z^{k!}=1\) 的 \(k!\) 个根都是 \(f(z)\) 的奇点，其中 \(k\) 为任意正整数。
3. 由此证明：不可能将 \(f(z)\) 延拓到单位圆外。单位圆周称为函数 \(f(z)\) 的自然边界。

\begin{solution}
1. 当 \(z=1\) 时，有 \(f(1)=\displaystyle\sum_{n=0}^\infty 1\) 发散，因此 \(1\) 为 \(f(z)\) 的奇点。

2. 当 \(z_0\) 为 \(z^{k!}=1\) 的根时，那么对于 \(n>k\)，有
   \[
   z_0^{n!} = z^{k!\cdot (k+1)\cdot (k+2)\cdots n} = 1^{(k+1)(k+2)\cdots n} = 1,
   \]
   因此有
   \[
   f(z_0) = \sum_{n=0}^{k-1} z_0^{n!} + \sum_{n=k}^\infty 1,
   \]
   级数发散，故 \(z_0\) 为 \(f(z)\) 的奇点。

3. 在单位圆上的奇点为 \(\ee^{2\pi \ii \cdot k/n!}\)，其中 \(n\in\mathbb{N},\,0\le k<n!\)，因此对于任意单位圆上的点，其邻域内一定存在奇点，因此不可能将 \(f(z)\) 延拓到单位圆外。
\end{solution}
\end{question}

\begin{question}
求证：
\[
\Gamma(z)=\int_L \ee^{-t} t^{z-1}\dd{t},\quad \operatorname{Re} z>0,
\]
\(L\) 是自原点出发的射线，\(0<|t|<\infty,\;|\arg t|<\pi/2\)。

\begin{solution}
\begin{center}
\begin{tikzpicture}
\draw[->] (0,0) -- (6,0) node[right] {$x$};
\draw[->] (0,-1) -- (0,4) node[above] {$y$};
\draw[thick,->] (0,0) -- (30:4) node[below right] {$L$};
\draw[thick,->] (0,0) -- (4,0) node[below] {$L'$};
\draw[thick] (4,0) arc (0:30:4);
\node at (4.5,1) {$C_R$};
\filldraw (0,0) circle (2pt);
\node at (0,0) [below left] {$O$};
\end{tikzpicture}
\end{center}
考察积分路径，由于 \(\ee^t t^{z-1}\) 在右半平面解析，因此由柯西积分定理
\[
\int_L \ee^{-t} t^{z-1}\dd{t}
= \int_{L'} \ee^{-t} t^{z-1}\dd{t} + \int_{C_R} \ee^{-t} t^{
\end{question}
\end{solution}
